\documentclass[12pt]{article}
\usepackage[utf8]{inputenc}
\usepackage[T1]{fontenc}
\usepackage{hyperref}
\title{Wampanoag Language Resources for K-12 Curriculum Development}
\author{Compiled by Claude Sonnet 5 (claude-sonnet-5) for Andr\'as Kornai}
\date{}
\begin{document}
\maketitle

\section{Introduction}
Wampanoag (W\^opan\^aak; ISO 639-3 wam; Glottocode wamp1249) is an Algonquian language of Massachusetts. It had no fluent speakers for generations and is being revived through the W\^opan\^aak Language Reclamation Project (WLRP), begun in 1993 under jessie ``little doe'' baird (Mashpee Wampanoag). A 2016 report counted 15 adult speakers and about 75 children at various levels of fluency. These figures may be out of date.

\section{Primary Educational Resources}
According to Cultural Survival, the WLRP runs the following programmes:
\begin{itemize}
\item Mukayuhsak Weekuw, a full-time Montessori immersion school for about 25 children from preschool to grade 1.
\item Weekly programming in public elementary schools.
\item World-language credit courses for grades 8--12 at the local middle and high school.
\item Elders', family and after-school classes, and immersion camps.
\end{itemize}
Cultural Survival profile: \url{https://www.culturalsurvival.org/node/12450}

No published lesson plans or curriculum documents were found in the results.

\section{Linguistic and Reference Materials}
The language was reconstructed from 17th-century sources, namely the Massachusett Bible, letters, wills and deeds. No dictionary or grammar URL was found in this search.

\section{Religious and Community Resources}
\begin{itemize}
\item Cultural Survival, further coverage of the project: \url{https://www.culturalsurvival.org/node/11598}
\end{itemize}

\section{Digital and Audio Resources}
\begin{itemize}
\item PRI story, ``Wampanoag language waking after nearly a century of sleep'' (2016): \url{https://www.pri.org/stories/2016/12/23/wampanoag-language-waking-after-nearly-century-sleep}
\end{itemize}

\section{Curriculum Development Context}
Unlike many languages in this survey, Wampanoag already has an active school-based revitalization effort led by the community, so the WLRP itself is the natural partner for any curriculum work. Its own website was not found in this search and should be located directly. Contact details are not given here.
\end{document}